\chapter{Celiac Disease}

\section*{Definition and Overview}
Celiac disease is a chronic autoimmune condition in which eating gluten -- a protein found in wheat, barley, and rye -- triggers the body's immune system to attack the lining of the small intestine. This damages the tiny, finger-like projections called villi that absorb nutrients, leading to malabsorption, digestive symptoms, and a range of effects across the body.

It is not a food allergy or an intolerance in the usual sense; it is a lifelong immune-mediated disease. The only reliable treatment is a strict, permanent gluten-free diet, which allows the intestine to heal and symptoms to resolve.

\section*{Epidemiology}
Celiac disease affects roughly 1 in 100 people worldwide, making it one of the most common chronic digestive conditions. It can present at any age, from early childhood to late adulthood, and appears slightly more often in females than in males.

The condition is more common in people of European descent and runs strongly in families: first-degree relatives have a significantly higher chance of being affected. Many cases are undiagnosed because symptoms are mild, varied, or absent, and awareness of the disease has grown substantially in recent decades, leading to more frequent testing and detection.

\section*{Aetiology and Pathophysiology}
Celiac disease develops through a combination of genetic susceptibility and environmental exposure to gluten:

\begin{itemize}
    \item \textbf{Genetic factors:} nearly all affected people carry the immune-system markers HLA-DQ2 or HLA-DQ8
    \item \textbf{Gluten exposure:} the gliadin component of gluten is the trigger in wheat, barley, and rye
    \item \textbf{Immune response:} abnormal immune cells react to gluten and damage the intestinal lining
    \item \textbf{Autoimmune nature:} the immune system also produces antibodies against tissue transglutaminase, an enzyme in the gut
    \item \textbf{Associated conditions:} type 1 diabetes, thyroid disease, and Down syndrome increase the risk
    \item \textbf{Environmental triggers:} events such as gut infections or surgery may sometimes coincide with the first onset of symptoms
\end{itemize}

The resulting inflammation flattens the villi and reduces the intestine's surface area, so the body struggles to absorb nutrients such as iron, folate, calcium, and vitamin B12.

\section*{Clinical Features}
Symptoms vary widely and may involve the gut, other body systems, or both:

\begin{itemize}
    \item Diarrhoea, loose stools, or pale, foul-smelling fatty stools
    \item Abdominal pain, bloating, and excess gas
    \item Weight loss or poor weight gain in children
    \item Fatigue and weakness, often from anaemia
    \item Iron-deficiency anaemia that does not respond to supplements
    \item Mouth ulcers and a blistering, intensely itchy skin rash (dermatitis herpetiformis)
    \item Bone or joint pain and low bone density
    \item In children: delayed growth, irritability, and failure to thrive
\end{itemize}

Some people, especially adults, have few or no digestive symptoms and are diagnosed after blood tests done for anaemia, fatigue, or other reasons.

\section*{Differential Diagnosis}
Several other conditions can cause similar symptoms and must be considered:

\begin{itemize}
    \item Irritable bowel syndrome (IBS)
    \item Non-celiac gluten sensitivity, in which symptoms improve on a gluten-free diet without intestinal damage
    \item Crohn's disease and other inflammatory bowel diseases
    \item Food allergies or lactose intolerance
    \item Tropical sprue and other causes of malabsorption
    \item Coeliac-like changes from some medicines or infections
    \item Chronic infections such as giardiasis
\end{itemize}

\section*{When to Seek Medical Care}
See a doctor if you or your child have persistent digestive symptoms, unexplained weight loss, chronic fatigue, or iron deficiency that does not respond to treatment.

\textbf{Call 000 (or local emergency services) for:}
\begin{itemize}
    \item Severe abdominal pain that is sudden or worsening
    \item Vomiting with inability to keep fluids down
    \item Significant dehydration or signs of severe illness
    \item Sudden collapse or anaphylactic-type symptoms, which are not typical of celiac disease and need urgent assessment
\end{itemize}

\section*{Diagnosis and Investigations}
\begin{itemize}
    \item \textbf{Blood tests:} screening for antibodies to tissue transglutaminase (tTG-IgA) and for total IgA level
    \item \textbf{Genetic testing:} HLA-DQ2/DQ8 typing can help rule out the disease
    \item \textbf{Endoscopy with biopsy:} small-intestine biopsy is the gold standard and shows the characteristic villous damage
    \item \textbf{Iron and nutrient studies:} checks for anaemia, folate, B12, calcium, and vitamin D deficiency
    \item \textbf{Bone density scan:} to assess for osteoporosis in newly diagnosed adults
    \item \textbf{Important:} continue eating gluten during testing, because a gluten-free diet can normalise blood tests and biopsy before diagnosis is made
\end{itemize}

\section*{Management}
\begin{itemize}
    \item \textbf{Strict gluten-free diet:} the cornerstone of treatment; excludes wheat, barley, rye, and their derivatives
    \item \textbf{Dietitian review:} to plan a nutritionally complete diet and identify hidden sources of gluten
    \item \textbf{Nutrient replacement:} supplements for iron, folate, B12, calcium, and vitamin D where deficient
    \item \textbf{Symptom control:} treatment of ongoing bloating or reflux that can persist despite diet
    \item \textbf{Dermatitis herpetiformis:} managed with a gluten-free diet, often with an initial course of dapsone
    \item \textbf{Follow-up:} repeat antibody testing to confirm diet response and monitor healing
    \item \textbf{Family screening:} testing of first-degree relatives
\end{itemize}

\section*{Complications}
\begin{itemize}
    \item Malnutrition and deficiency of iron, folate, B12, and fat-soluble vitamins
    \item Osteoporosis and increased fracture risk from poor calcium absorption
    \item Anaemia and chronic fatigue
    \item Reduced fertility in some untreated cases
    \item Delayed growth and development in children
    \item Refractory celiac disease and a small increased risk of lymphoma in people with longstanding untreated disease
    \item Increased risk of other autoimmune conditions
\end{itemize}

\section*{Prognosis and Outcomes}
On a strict gluten-free diet, most people feel better within weeks to months, and the intestinal lining usually heals over one to two years. Many people have a normal life expectancy and few ongoing problems. Some continue to have symptoms from associated conditions such as IBS or lactose intolerance. Outcomes are best with early diagnosis and consistent dietary adherence, while untreated disease carries higher risks of complications such as osteoporosis and malnutrition.

\section*{Prevention and Self-Care}
\begin{itemize}
    \item Follow the gluten-free diet strictly and permanently; even small amounts cause damage
    \item Learn to read food labels and watch for hidden gluten in sauces, processed foods, and medications
    \item Avoid cross-contamination in shared kitchens and when eating out
    \item Work with a dietitian to keep the diet balanced and varied
    \item Take prescribed supplements to correct deficiencies
    \item Have regular check-ups, including repeat antibody tests and bone density scans where recommended
    \item Join a celiac support group for practical advice and meal planning help
\end{itemize}

\section*{Sources and Further Reading}
The following organisations provide reliable information on celiac disease and gluten-free living.

\begin{itemize}
    \item National Health Service. \textit{NHS conditions guide on coeliac disease.} https://www.nhs.uk/conditions/
    \item Mayo Clinic. \textit{Celiac disease: symptoms, causes, and treatment.} https://www.mayoclinic.org/
    \item World Health Organization. \textit{Global guidance on digestive and nutritional health.} https://www.who.int/
    \item MedlinePlus. \textit{Consumer health information on celiac disease.} https://medlineplus.gov/
    \item Centers for Disease Control and Prevention. \textit{Nutrition and digestive health information.} https://www.cdc.gov/
    \item MSD Manual. \textit{Celiac disease -- professional overview.} https://www.msdmanuals.com/
\end{itemize}
